\chapter{Título do capitulo 2}\label{chap2}
    

\section{Exemplo de uso de label, ref e todo}
\label{sec:introducao}

Este é um exemplo simples para mostrar como usar \texttt{\textbackslash label},
\texttt{\textbackslash ref}, \texttt{\textbackslash todo}, \texttt{\textbackslash todo[inline]},
\texttt{refcheck} e \texttt{showkeys}.

% todo na margem
\todo{Revisar}

Na Seção \ref{sec:introducao}, começamos a explicar o uso de referências internas.

\begin{equation}
\label{eq:principal}
x^2+y^2=1
\end{equation}

A equação \eqref{eq:principal} representa a circunferência unitária.

% todo inline dentro do texto
\todo[inline]{Melhorar a explicação geométrica da equação \eqref{eq:principal}.}

\subsection{Mais um exemplo}
\label{subsec:exemplo}

Agora fazemos referência novamente à Seção \ref{sec:introducao}
e à equação \eqref{eq:principal}.

\begin{teorema}
\label{thm:teste}
Se $x=y=0$, então a equação \eqref{eq:principal} não é satisfeita.
\end{teorema}

O Teorema \ref{thm:teste} é apenas ilustrativo.


% Comece aqui!sss